
\subsubsection{Editorial-Oriented Two-Stage Paradigm}

A preliminary analysis is conducted using a standard linear text classifier based on TF--IDF representations (word- and character-level n-grams), source indicators, and simple numeric features. Despite extensive hyperparameter exploration, performance consistently converges to a narrow plateau, with optimal configurations characterized by:
\begin{table}[H]
\centering
\caption{Optimal hyperparameter ranges for the linear TF--IDF baseline.}
\label{tab:baseline_hyperparams}
\begin{tabular}{lc}
\hline
\textbf{Hyperparameter} & \textbf{Selected value / range} \\
\hline
Word n-grams        & 2 \\
Character n-grams   & 5--7 \\
Regularization $C$  & 0.6--1.0 \\
Max document freq.  & 0.85--0.9 \\
\hline
\end{tabular}
\end{table}


Across these settings, the macro-averaged F1 score saturates around $0.719$, indicating that a single-stage linear model has largely exhausted the information available under a homogeneous formulation.

\begin{table}[H]
\centering
\caption{Examples of high-support editorial tokens with near-zero conditional entropy.}
\label{tab:low_entropy_tokens}
\small
\resizebox{\linewidth}{!}{%
\begin{tabular}{lcccl}
\hline
\textbf{Token} & \textbf{Type} & \textbf{Support} & \textbf{Entropy} & \textbf{Dominant label} \\
\hline
\texttt{<h4>} & editorial & 628 & 0.00 & 2 \\
\texttt{</h4>} & editorial & 628 & 0.00 & 2 \\
\texttt{<p align="right">} & editorial & 311 & 0.00 & 2 \\
\texttt{<a href=...>} & editorial & 233 & 0.00 & 2 \\
\texttt{<img class="rss-ad">} & editorial & 196 & 0.00 & 2 \\
\texttt{<strong>} & editorial & 111 & 0.00 & 2 \\
\texttt{<cite>} & editorial & 40 & 0.00 & 2 \\
\texttt{<em>} & editorial & 57 & 0.47 & 2 \\
\hline
\end{tabular}}
\end{table}

Inspection of the learned feature weights reveals that the most influential predictors are predominantly editorial and structural rather than semantic, including publisher identifiers, RSS-related markers, and HTML or URL fragments. These features exhibit strong class-specific associations and dominate the linear decision function.

To formalize this behavior, we analyze token-level statistics. Let $t$ denote a token and $Y$ the class label. The empirical conditional entropy is defined as:
\[
H(Y \mid t) = - \sum_{c \in \mathcal{Y}} p(c \mid t)\log p(c \mid t).
\]
For a non-trivial subset of tokens, we observe $H(Y \mid t)=0$, meaning that their presence alone determines the class on the development set with certainty. This evidences a structural heterogeneity in the dataset, motivating an explicit decomposition of the classification process.

We therefore adopt a Two-Stage classification paradigm. For each token $t$, we compute:
\[
\text{purity}(t) = \max_{c \in \mathcal{Y}} p(c \mid t), \quad
\text{support}(t) = |\{x_i : t \in x_i\}|.
\]
Based on these quantities, we define a rule set:
\[
\mathcal{R} = \{t \in \mathcal{T} : \text{purity}(t) \ge \tau_p \;\wedge\; \text{support}(t) \ge \tau_s\},
\]
where $\tau_p$ and $\tau_s$ are fixed thresholds.

\paragraph{Stage 1: Deterministic Assignment}
Given a document $x$, if there exists a token $t \in \mathcal{R}$ such that $t \in x$, the document is assigned the label:
\[
\hat{y}^{(1)}(x) = \arg\max_{c \in \mathcal{Y}} p(c \mid t).
\]

\paragraph{Stage 2: Probabilistic Classification}
Documents not covered by Stage~1 are processed by a probabilistic classifier $f_\theta(x)$ operating on TF--IDF representations and auxiliary features:
\[
\hat{y}^{(2)}(x) = \arg\max_{c \in \mathcal{Y}} f_\theta(x)_c.
\]

\paragraph{Unified Decision Rule}
The final prediction is given by:
\[
\hat{y}(x) =
\begin{cases}
\hat{y}^{(1)}(x), & \text{if } \exists\, t \in \mathcal{R} \text{ such that } t \in x, \\
\hat{y}^{(2)}(x), & \text{otherwise}.
\end{cases}
\]

This formulation explicitly separates deterministic editorial patterns from ambiguous instances requiring probabilistic reasoning. In the absence of high-purity patterns, the model degrades gracefully to the second stage, recovering the standard single-stage classifier without loss of generality.
